\section*{Abstract}

\subsection{Objectives} This study synthesizes evidence on outcomes following intentional ingestion of foreign objects (IIFO) and assesses how patient motivation, object characteristics, and demographics affect the need for intervention, complications, and mortality.

\subsection{Methods} We conducted a systematic search across multiple databases, including PubMed, Embase, and Scopus, for studies on the intentional ingestion of non-nutritive foreign bodies in humans. Eligible studies reported on endoscopic, surgical, or conservative outcomes, excluding accidental ingestions, non-English reports, and animal studies. A single reviewer performed initial screenings, with 10\% double-checked for accuracy. We extracted data at both study and case levels, utilizing meta-analysis and $\chi^2$ testing for statistical analysis.

\subsection{Results} We included 71 case reports and three case series. Pooled outcome rates revealed 50\% underwent endoscopy, 30\% surgery, and 20\% conservative management. Complications occurred in 30\% of cases, with mortality under 1\%. The predominant motivations for ingestion were protest (approximately 80\%) and involved sharp objects (approximately 90\%). Protest motivation lowered surgical intervention rates, while intent-to-harm increased surgical risks. Adults aged 40--64 exhibited lower surgical rates but higher mortality.

\subsection{Conclusions} Patient motivation significantly impacts IIFO outcomes, necessitating improved reporting practices in future studies.